% ID: SEL-D0082-OA-E03
% ESTADO: cotejada. Segunda pasada visual y borrador compilado revisados.
% Fuente DOC: Selectividad/sel_2012_matematicas_aplicadas/4 MATEMÁT. APLIC. CC. SS/EXAMEN 4-ANDALUCÍA.doc
% SHA-256 original: 8446c5da6b4b7cfc0f48b89d1b2d052ae0c8a1c4b5bc0b86b7e63399ce8656f8
% Página convertida: 1; opción A; ejercicio 3.

Una compañía de seguros ha hecho un seguimiento durante un año a 50000 coches de la marca A, a 20000 de la marca B y a 30000 de la C, que tenía asegurados, obteniendo que, de ellos, habían tenido accidente 650 coches de la marca A, 200 de la B y 150 de la C. A la vista de estos datos:
\begin{enumerate}[label=\alph*)]
\item \textbf{(1.25 puntos)} \textquestiondown{}Cuál de las tres marcas de coches tiene menos proporción de accidentes?
\item \textbf{(1.25 puntos)} Si, elegido al azar uno de los coches observados, ha tenido un accidente, \textquestiondown{}cuál es la probabilidad de que sea de la marca C?
\end{enumerate}
